\chapter{Machine Learning Model Development}
\label{ch:ml-development}

\section{Dataset and Splits}

The feature matrix contains 5,257 daily records spanning May~2011 to
December~2025. A strict chronological split was applied to prevent data leakage:

\begin{table}[H]
\centering
\caption{Train / Validation / Test Split}
\begin{tabular}{llr}
\toprule
\textbf{Split} & \textbf{Period} & \textbf{Rows} \\
\midrule
Train      & 2011-05-18 → 2022-12-31 & 4,160 \\
Validation & 2023-01-01 → 2023-12-31 &   365 \\
Test       & 2024-01-01 → 2025-12-31 &   731 \\
\bottomrule
\end{tabular}
\end{table}

All preprocessing (median imputation, scaling) was fitted exclusively on the
training set and applied to validation and test sets without look-ahead.

\section{Baseline}

A persistence baseline was established using yesterday's PM2.5 value as the
forecast for all horizons. This represents the simplest possible model and
provides a floor that any learned model must surpass.

\begin{table}[H]
\centering
\caption{Persistence Baseline — Test MAE}
\begin{tabular}{lr}
\toprule
\textbf{Horizon} & \textbf{MAE (µg/m³)} \\
\midrule
t+1 & 5.30 \\
t+3 & 8.97 \\
t+7 & 10.93 \\
\bottomrule
\end{tabular}
\end{table}

\section{Model Selection}

Three model families were evaluated under an Optuna-driven hyperparameter search
(100+ trials per family per horizon), all sharing the same 45-feature input and
chronological split:

\textbf{LightGBM} was selected as the \emph{champion}. It combines the
efficiency of histogram-based boosting with strong tabular performance and sub-5ms
inference per row, and outperformed both competitors on every horizon. Best
hyperparameters at t+1 include 357 estimators, 255 leaves, max\_depth=6,
learning\_rate $\approx$ 0.043; longer horizons use deeper, more regularised
configurations.

\textbf{XGBoost} was retained as a \emph{strong secondary} and runtime fallback:
it is architecturally close to the LightGBM champion, trails it by $<$~0.3~MAE
at every horizon, and is exposed on the FastAPI service as
\texttt{?model=xgboost} for A/B comparison.

\textbf{LSTM} was trained as a \emph{comparison baseline} to assess whether
temporal sequence structure adds value over tabular lag features. The tuned
architecture (Optuna search over arch~$\in$~\{LSTM, Attention, CNN-LSTM\},
seq\_len, hidden units, dropout, learning rate) converges to architecture-horizon
pairs (attention @ seq=7 for t+1, CNN-LSTM @ seq=30 for t+3).

\section{Results}

\subsection{LightGBM (Champion)}

\begin{table}[H]
\centering
\caption{LightGBM Test Set Performance}
\label{tab:lgbm_results}
\begin{tabular}{lrrrr}
\toprule
\textbf{Horizon} & \textbf{MAE} & \textbf{RMSE} & \textbf{R²} & \textbf{AUROC} \\
\midrule
t+1 & 5.12 & 9.00  & 0.841 & 0.992 \\
t+3 & 6.77 & 11.78 & 0.729 & 0.973 \\
t+7 & 8.25 & 14.36 & 0.600 & 0.946 \\
\bottomrule
\end{tabular}
\end{table}

LightGBM beats the persistence baseline at every horizon: MAE improvement is
+3\% at t+1, +25\% at t+3, and +25\% at t+7. The +3\% lift at t+1 is modest
because same-day PM2.5 is an extremely strong anchor (the lag-1 signal), and
any learned model must push past that ceiling using fire + weather features.

\subsection{XGBoost (Strong Secondary)}

\begin{table}[H]
\centering
\caption{XGBoost Test Set Performance (Optuna-tuned)}
\label{tab:xgb_results}
\begin{tabular}{lrrrr}
\toprule
\textbf{Horizon} & \textbf{MAE} & \textbf{RMSE} & \textbf{R²} & \textbf{AUROC} \\
\midrule
t+1 & 5.31 & 8.95  & 0.843 & 0.992 \\
t+3 & 7.08 & 11.90 & 0.723 & 0.973 \\
t+7 & 8.90 & 15.32 & 0.544 & 0.933 \\
\bottomrule
\end{tabular}
\end{table}

XGBoost trails the LightGBM champion by 0.19--0.65~MAE across horizons. The gap
widens at t+7, where LightGBM's histogram-based leaf growth generalises slightly
better on the harder long-range task. XGBoost is kept loaded in the FastAPI
service as a runtime fallback (\texttt{?model=xgboost}) and as an independent
score for threshold-based A/B comparison.

\subsection{LSTM (Comparison)}

\begin{table}[H]
\centering
\caption{LSTM Test Set Performance}
\label{tab:lstm_results}
\begin{tabular}{lrrr}
\toprule
\textbf{Horizon} & \textbf{MAE} & \textbf{RMSE} & \textbf{R²} \\
\midrule
t+1 & 6.67 & 11.64 & 0.753 \\
t+3 & 8.87 & 15.89 & 0.540 \\
t+7 & 8.06 & 14.49 & 0.608 \\
\bottomrule
\end{tabular}
\end{table}

Both gradient-boosted models outperform LSTM at t+1 (MAE 5.12 / 5.31 vs.\ 6.67)
and t+3 (6.77 / 7.08 vs.\ 8.87), while performance at t+7 is comparable
(8.25 / 8.90 vs.\ 8.06). The LSTM's weaker t+3 result (R²=0.54) reflects the
difficulty of the mid-range horizon, where the recency of lag features decays
but seasonal signals are not yet dominant. LightGBM is selected as the champion.

\subsection{Actual vs.\ Predicted}

\begin{figure}[H]
    \centering
    \includegraphics[width=0.9\textwidth]{fig_actual_vs_predicted}
    \caption{Champion t+1 actual vs.\ predicted PM2.5 on the 2024--2025 test set.
             The model captures seasonal haze peaks accurately.}
    \label{fig:actual_vs_predicted}
\end{figure}

\section{Ablation Study: FIRMS Fire Data}

To quantify the contribution of NASA FIRMS fire hotspot features, Scenario~D
retrained the XGBoost candidate with the six FIRMS-derived features removed
(39 features total) and compared against the full-feature XGBoost run from the
same training harness. Using XGBoost rather than the LightGBM champion isolates
the ablation from concurrent hyperparameter differences.

\begin{table}[H]
\centering
\caption{Scenario D — MAE Impact of Removing FIRMS Features (XGBoost)}
\begin{tabular}{lrrr}
\toprule
\textbf{Horizon} & \textbf{Full MAE} & \textbf{No-FIRMS MAE} & \textbf{MAE increase} \\
\midrule
t+1 & 6.11 & 6.52 & +6.8\% \\
t+3 & 7.37 & 7.75 & +5.1\% \\
t+7 & 8.44 & 8.80 & +4.3\% \\
\bottomrule
\end{tabular}
\end{table}

Fire data provides consistent, measurable improvement across all horizons.
The effect is strongest at t+1 (6.8\%), where recent hotspot counts directly
predict tomorrow's PM2.5 accumulation.

\begin{figure}[H]
    \centering
    \includegraphics[width=0.75\textwidth]{fig_firms_ablation}
    \caption{MAE comparison: full model vs.\ model without FIRMS fire features.}
    \label{fig:firms_ablation}
\end{figure}

\section{Alert Classifier Analysis (Scenario C)}

The binary hazard alert (PM2.5 $>$ 50~\textmu g/m\textsuperscript{3}) is derived
from the regression output by applying a threshold. Scenario~C investigated
whether the default 50~\textmu g/m\textsuperscript{3} threshold is optimal by
computing F1 score across all possible thresholds on the test set.

\begin{table}[H]
\centering
\caption{Scenario C — Alert Threshold Optimisation (XGBoost scores)}
\begin{tabular}{lrrrr}
\toprule
\textbf{Horizon} & \textbf{Default (50)} & \textbf{Optimal threshold} & \textbf{Default F1} & \textbf{Optimal F1} \\
\midrule
t+1 & 50.0 µg/m³ & 53.7 µg/m³ & 0.839 & \textbf{0.879} \\
t+3 & 50.0 µg/m³ & 51.0 µg/m³ & 0.787 & \textbf{0.789} \\
t+7 & 50.0 µg/m³ & 45.8 µg/m³ & 0.684 & \textbf{0.695} \\
\bottomrule
\end{tabular}
\end{table}

Threshold tuning was performed against the XGBoost test-set predictions (the
regression scores the same pipeline exposes as \texttt{?model=xgboost}) so that
each alert class is evaluated on an unbiased held-out split. The standard
50~\textmu g/m\textsuperscript{3} threshold is near-optimal for t+3 and t+7.
For t+1, raising the threshold to 53.7~\textmu g/m\textsuperscript{3} reduces
false alarms (precision 0.769 → 0.885) while maintaining recall
(0.924 → 0.873), yielding a +4.0\% F1 gain.

\begin{figure}[H]
    \centering
    \includegraphics[width=\textwidth]{fig_pr_curve_t1}
    \caption{Precision-recall curve and F1 vs.\ threshold for t+1 alert classifier.}
    \label{fig:pr_curve_t1}
\end{figure}

\section{Model Interpretation (SHAP)}

SHAP (SHapley Additive exPlanations) values were computed against the XGBoost
strong-secondary on the test set. XGBoost shares LightGBM's underlying gradient-
boosted tree structure and exposes the same 45-feature input, so its SHAP
attributions are representative of the broader gradient-boosting regime; running
SHAP on XGBoost avoids LightGBM's TreeExplainer compatibility caveats while
preserving the physical interpretation story.

\begin{figure}[H]
    \centering
    \includegraphics[width=0.85\textwidth]{fig_shap_summary_t1}
    \caption{SHAP summary plot at t+1. \texttt{pm25\_lag1} dominates,
    confirming that next-day PM2.5 is primarily driven by today's concentration.}
    \label{fig:shap_t1}
\end{figure}

\begin{figure}[H]
    \centering
    \includegraphics[width=0.85\textwidth]{fig_shap_summary_t3}
    \caption{SHAP summary plot at t+3. \texttt{hotspot\_14d\_roll} becomes
    the dominant signal, reflecting fire accumulation over two weeks.}
    \label{fig:shap_t3}
\end{figure}

A key insight is that the dominant predictive signal shifts with horizon:
\begin{itemize}
    \item \textbf{t+1:} Recent PM2.5 (\texttt{pm25\_lag1}, mean SHAP = 6.85) ---
          next-day concentration is primarily a persistence phenomenon.
    \item \textbf{t+3:} 14-day fire accumulation (\texttt{hotspot\_14d\_roll}, 4.63)
          --- two-week fire buildup is the strongest medium-range signal.
    \item \textbf{t+7:} Seasonal position (\texttt{sin\_month}, 4.63;
          \texttt{hotspot\_7d\_roll}, 3.00) --- at 7 days, the model relies on
          calendar position within the haze season plus near-term fire activity.
\end{itemize}

This gradient from measurement-driven to climatology-driven prediction is
consistent with the physical understanding of PM2.5 dynamics in Chiang Mai.
